\documentclass[12pt,letterpaper]{article}

% Arial (Helvetica como equivalente portable para pdflatex)
\usepackage[scaled=1.0]{helvet}
\renewcommand{\familydefault}{\sfdefault}
\usepackage[T1]{fontenc}
\usepackage[utf8]{inputenc}

% Idioma español (Mexico)
\usepackage[spanish,mexico]{babel}

% Márgenes — tamaño carta, 2.5 cm todos los lados
\usepackage[letterpaper, margin=2.5cm]{geometry}

% Interlineado sencillo
\usepackage{setspace}
\singlespacing
\setlength{\parskip}{4pt}
\setlength{\parindent}{0pt}

% Secciones compactas
\usepackage{titlesec}
\titlespacing*{\section}{0pt}{8pt}{3pt}

% Hipervínculos sin color visible
\usepackage[hidelinks]{hyperref}

% Unidades SI
\usepackage{siunitx}

\begin{document}

% ============================================================
% Encabezado
% ============================================================
\begin{center}
    {\large\textbf{Sistema de teleoperación de robot móvil-manipulador\\
    para logística inclusiva}}\\[5pt]
    Sebastián Méndez Villegas$^{1}$, Santiago Cuesta Machuca$^{1}$\\[2pt]
    {\small $^{1}$Departamento de Estudios en Ingeniería para la Innovación,
    Universidad Iberoamericana Ciudad de México}\\[1pt]
    {\small \texttt{a2229332@correo.uia.mx}}
\end{center}
\vspace{4pt}\hrule\vspace{6pt}

% ============================================================
\section{Introducción}
% ============================================================

En México, las personas con discapacidad motriz en miembros inferiores enfrentan
tasas de desocupación significativamente superiores al promedio nacional y acceso
limitado a sectores que demandan movilidad física continua, como la logística y la
distribución de mercancía \cite{inegi2023}. La teleoperación robótica mediada por
interfaces de realidad extendida (XR) representa una alternativa tecnológica viable
para habilitar trabajo remoto de alta complejidad desde posiciones estáticas, reduciendo
barreras de inclusión laboral \cite{sheridan1992}. Investigaciones recientes han
demostrado la utilidad de entornos de realidad virtual para aprender y ejecutar tareas
de teleoperación \cite{lipton2018}, motivando el desarrollo de sistemas accesibles y
de bajo costo. Este trabajo presenta el diseño, construcción e integración de un
sistema completo de teleoperación compuesto por un robot móvil diferencial, un
manipulador de tres grados de libertad (GDL) con pinza, y una interfaz inmersiva en
Meta Quest~3.

% ============================================================
\section{Objetivo}
% ============================================================

Diseñar, construir y validar un sistema de teleoperación modular compuesto por una
plataforma AGV diferencial con manipulador de 3~GDL y pinza de agarre, controlado
mediante una interfaz de realidad mixta sobre Meta Quest~3, orientado a facilitar
tareas logísticas remotas para personas con discapacidad motriz en miembros inferiores.

% ============================================================
\section{Métodos}
% ============================================================

El sistema se desarrolló en tres subsistemas integrados: plataforma robótica,
electrónica-firmware embebido, y middleware de coordinación con interfaz XR.

\textbf{Plataforma robótica.}
La base móvil es un AGV diferencial rehabilitado con estructura de acero calibre~18,
dos motores DC con encoders y cuatro ruedas locas (castor wheels) \cite{siegwart2011}.
El manipulador de 3~GDL fue diseñado completamente en Autodesk Inventor y fabricado
con acero cortado a láser (calibres 0.6 y 1.2~mm); los actuadores son motores NEMA17
con reductor planetario 10:1 y transmisión por banda HTD3M relación 2.5:1 (reducción
compuesta 25:1) \cite{siciliano2010}. La pinza de tipo piñón-cremallera (impresión 3D)
opera con motor Pololu 100:1 y recorrido de \SI{80}{mm} controlado mediante encoder de
cuadratura.

\textbf{Electrónica embebida.}
Se diseñó una PCB propia denominada \textit{Módulo Puente H ESP32-C3} en KiCad
(12--24~V, 4~A nominal, 8~A pico, MOSFETs IRF540N/IRF9540N con aislamiento óptico 4N25).
Tres microcontroladores ESP32-C3 ejecutan firmware en MicroPython: uno como maestro I\textsuperscript{2}C
para control diferencial de tracción, uno como esclavo I\textsuperscript{2}C para
generación de señales PUL/DIR hacia los drivers CL57T del manipulador, y uno por
USB-CDC para control de posición del gripper con perfil trapezoidal de velocidad y
detección de atasco.

\textbf{Middleware y percepción.}
Una computadora Intel NUC~i7 embarcada coordina el sistema con Python~3.12 y ZeroMQ.
Los sensores incluyen una Intel RealSense~D435i (RGB-D, 640$\times$480 a 30~fps) con
detección de objetos mediante YOLOv8~nano, y un RPLiDAR~C1 (LiDAR 2D, 12~Hz) para
generación de grillas de ocupación en tres escalas (1, 2 y 3~m de radio, resolución
de celda 1~cm).

\textbf{Interfaz XR.}
Una aplicación Unity compilada para Meta Quest~3 (Meta XR SDK~83.0) presenta en
passthrough: panel de video RGB, grilla de ocupación LiDAR en 2D, controles
deslizantes para los 3~GDL del manipulador, y panel de telemetría del sistema
(posición de articulaciones, posición del gripper, estado de sensores), con
retroalimentación háptica y sonora por los controladores del headset.

% ============================================================
\section{Resultados}
% ============================================================

La PCB fue fabricada, ensamblada y validada con temperatura máxima medida en los
MOSFETs de \SI{61}{\degreeCelsius} bajo carga continua a \SI{60}{\minute} con
disipadores pasivos, dentro del rango operativo admisible. El manipulador alcanzó
todos los rangos de movimiento definidos: base $\pm$80°, codo 0°--136.5°, muñeca
$-$220°--0°, con resoluciones de 27.78 y 69.44~pasos/grado respectivamente, sin
pérdida de pasos en movimientos estáticos. La pinza completó el recorrido de
\SI{80}{mm} sin pérdida de posición y activó la detección de obstrucción mecánica
(stall) en menos de \SI{500}{\milli\second}; la calibración de extremos se recupera
automáticamente tras un reinicio desde memoria JSON. El streaming de video fue
validado a 30~fps y el ciclo de comandos de tracción opera a 15~Hz con latencia
total inferior a \SI{50}{\milli\second}. La interfaz XR integra en tiempo real el
video, la grilla LiDAR en sus tres modos, el control 3D del manipulador y la pinza
desde el Meta Quest~3. Se verificó la teleoperación completa de extremo a
extremo ---tracción, manipulación de 3~GDL y agarre--- con realimentación visual y
háptica; el watchdog ZeroMQ detiene la tracción en \SI{350}{\milli\second} ante
pérdida de enlace.

% ============================================================
\section{Conclusiones}
% ============================================================

El proyecto demuestra la viabilidad técnica de un sistema de teleoperación inmersivo
para logística inclusiva construido íntegramente con hardware diseñado a medida y
software de código abierto. La arquitectura modular basada en ZeroMQ, ESP32-C3 y
Unity facilita la integración de subsistemas heterogéneos con baja latencia y alta
cohesión funcional. La validación formal con usuarios objetivo ---personas con
discapacidad motriz en miembros inferiores--- constituye el paso inmediato para
cuantificar el impacto del sistema en empleabilidad y desempeño operativo. Como
trabajo futuro se contemplan la incorporación de SLAM para mapeo autónomo del entorno,
la transmisión RGB-D completa al headset y la integración de retroalimentación de
fuerza en la pinza.

% ============================================================
\section*{Agradecimientos}
% ============================================================

Los autores agradecen el apoyo de la Universidad Iberoamericana Ciudad de México por
los recursos y espacios de laboratorio utilizados en el desarrollo de este proyecto.

% ============================================================
\begin{thebibliography}{9}
\addcontentsline{toc}{section}{Bibliografía}

\bibitem{inegi2023}
INEGI (2023). \textit{Encuesta Nacional sobre Discriminación 2022 (ENADIS)}.
Instituto Nacional de Estadística y Geografía, México.
\url{https://www.inegi.org.mx/programas/enadis/2022/}

\bibitem{sheridan1992}
Sheridan, T.~B. (1992).
\textit{Telerobotics, Automation, and Human Supervisory Control}.
MIT Press.

\bibitem{lipton2018}
Lipton, J.~I., Fay, A.~J., \& Rus, D. (2018).
Baxter's Homunculus: Virtual Reality Spaces for Teleoperation Learning.
\textit{IEEE Robotics and Automation Letters}, 3(3), 2344--2351.

\bibitem{siciliano2010}
Siciliano, B., Sciavicco, L., Villani, L., \& Oriolo, G. (2010).
\textit{Robotics: Modelling, Planning and Control}.
Springer-Verlag.

\bibitem{siegwart2011}
Siegwart, R., Nourbakhsh, I.~R., \& Scaramuzza, D. (2011).
\textit{Introduction to Autonomous Mobile Robots}, 2nd~ed.
MIT Press.

\end{thebibliography}

\end{document}
